\textbf{Speech Reconstruction}~
\begin{table*}[t]
    \setlength{\belowcaptionskip}{-0.2cm}
    % \renewcommand{\arraystretch}{1.2}
    \Large
    \centering
    \resizebox{0.5\textwidth}{!}{
    \begin{tabular}{l|cc|c}
        \toprule
        ~ & \multicolumn{2}{c}{\textbf{Objective}} & \textbf{Subjective}  \\
        Tokenizer & WER$\downarrow$ & VISQOL$\uparrow$ & MUSHRA$\uparrow$ \\
        \midrule
        Groundtruth & 4.58 & - & 91.46 \\
        \midrule
        EnCodec & 5.11 & 4.37 &  79.86\\
        \method & 5.04 & 4.30 &  90.55\\
        \bottomrule
    \end{tabular}}
    \caption{Results of speech reconstruction
    }
    \label{tab:speech_reconstruction}
    % \vspace{-0.3em}
\end{table*}
Table~\ref{tab:speech_reconstruction} summarizes the results of speech reconstruction experiments. The \method achienves lower WER than Encodec, demonstrating its superior ability to preserve content. Additionally, \method attains a comparable VISQOL score but a higher MUSHRA score than EnCodec, which indicates its stronger capability in generating high-quality speech. 
% We put samples of speech reconstruction in our demo page~\footnote{\url{https://0nutation.github.io/SpeechTokenizer.github.io/}}.


% \begin{table*}[t]
%     % \setlength{\abovecaptionskip}{-0.cm}
%     \setlength{\belowcaptionskip}{-0.2cm}
%     % \renewcommand{\arraystretch}{1.2}
%     \Large
%     \centering
%     \resizebox{0.9\textwidth}{!}{
%     \begin{tabular}{l|ccc|ccc}
%         \toprule
%         ~ & \multicolumn{3}{c}{\textbf{Content}} & \multicolumn{3}{c}{\textbf{Timbre}}  \\
%         ~ & \multicolumn{2}{c}{DSRIBench} & Reconstruction &
%         \multicolumn{2}{c}{DSRIBench} & Reconstruction \\
%         ~ & MI & WER & WER& MI & ACC & SIM \\
%         \midrule
%         EnCodec-L1 & 16.5 & 61.52 & 38.34 & 6.52 & 0.44 & 0.92 \\
%         EnCodec-All & 23.6 & 30.91 & 5.11 & 14.7 & 0.94 & 0.98 \\
%         \method-L1 & & & 9.57 & & & 0.74 \\
%         \method-L2:8 & & & 86.58 & & & 0.79 \\
%         \method-all & & & 5.04 & & & 0.97 \\
%         \bottomrule
%     \end{tabular}}
%     \caption{Results of speech disentanglement.
%     }
%     \label{tab:speech_disentanglement}
% \end{table*}


\begin{table*}[t]
    % \setlength{\abovecaptionskip}{-0.cm}
    \setlength{\belowcaptionskip}{-0.2cm}
    % \renewcommand{\arraystretch}{1.2}
    \Large
    \centering
    \resizebox{0.9\textwidth}{!}{
    \begin{tabular}{lcc|cc|cc}
        \toprule
        ~ & & & \multicolumn{2}{c}{\textbf{Text Alignment}} & \multicolumn{2}{c}{\textbf{Information Preservation}}   \\
        Tokenizer & & Teacher & MI$\uparrow$ & WER$^{\dagger}\downarrow$ & WER$^{\ast}\downarrow$ & SIM$\uparrow$\\
        \midrule
        Groundtruth & & & - & - & 4.58 & 1.0 \\
        HuBERT & KM500 & - & 31.2 & 9.88 &16.26 & 0.77\\
        EnCodec & RVQ-1 & - & 16.5 & 61.52 & 38.34 & 0.92 \\
        EnCodec & RVQ-1:8 & - & 23.6 & 30.91 & 5.11 & 0.98 \\
        \midrule
        \textit{Ablations} & & & & & & \\
        \method & RVQ-1 & HuBERT avg  & 30.9 & 15.58 & 9.57 & 0.74 \\
        \method & RVQ-1:8 & HuBERT avg &29.7  & 16.03 & 5.04 & 0.97 \\
        \method & RVQ-1 & HuBERT L9 & 32.9 & 12.68 & 14.17&0.73  \\
        \method & RVQ-1:8 & HuBERT L9 & 31.6 & 13.12 & 5.31 & 0.97 \\
        \method & RVQ-1 & HuBERT units & 24.2 & 34.13 & 20.02 & 0.72 \\
        \method & RVQ-1:8 & HuBERT units & 25.1 & 30.71 & 5.84 & 0.95 \\
        
        
        \bottomrule
    \end{tabular}}
    \caption{Results on \bench. MI and WER$^{\dagger}$ refer to mutual information and word error rate of the downstream model. WER$^{\ast}$ and SIM refer to word error rate and speaker similarity of resynthesized speech respectively. RVQ-$n$ denotes the tokens of the $n^{th}$ RVQ layer. RVQ-$n$:$m$ denotes the tokens from the $n^{th}$ layer to the $m^{th}$ layer.
    }
    \label{tab:bench_results}
\vspace{-0.1em}
\end{table*}

\textbf{Performance on \bench }~
Table~\ref{tab:bench_results} displays the performance of \method on \bench. 
Compared with EnCodec-RVQ-1, \method-RVQ-1 achieves higher mutual information between text and lower WER of downstream model. This suggests that \method exhibits a stronger alignment with textual content. Meanwhile, the of resynthesized speech of \method RVQ-1 tokens achieves lower WER and speaker similarity, indicating its capability to retain more content-related information while disregarding timbre characteristics, similar to semantic tokens.
The resynthesized speech of \method RVQ-1:8 tokens demonstrates low WER and high speaker similarity, illustrating \method's competence in preserving comprehensive speech information, similar to acoustic tokens. 
Furthermore, the speaker similarity of resynthesized speech of \method RVQ-1 tokens is notably low, whereas that of \method RVQ-1:8 tokens is considerably high. This observation implies that the tokens from subsequent layers compensate for the timbre information that is discarded by the first layer tokens.




\textbf{Zero-shot TTS}~
As shown in Table~\ref{tab:zstts}, our USLM demonstrates lower WER than VALL-E. This result highlights that \method can contribute to a more precise modeling of content information. Additionally, the USLM demonstrates superior speaker similarity, implying that a decoupled information structure is more conducive to modeling speaker-related information. 
% We put samples of zero-shot TTS on our demo page~\footnote{\url{https://0nutation.github.io/SpeechTokenizer.github.io/}}.


\begin{table*}[t]
    \setlength{\belowcaptionskip}{-0.2cm}
    % \renewcommand{\arraystretch}{1.2}
    \Large
    \centering
    \resizebox{0.65\textwidth}{!}{
    \begin{tabular}{lc|cc|cc}
        \toprule
        ~ & & \multicolumn{2}{c}{\textbf{Objective}} & \multicolumn{2}{c}{\textbf{Subjective}}  \\
        Model & Tokenizer & WER$\downarrow$ & SIM$\uparrow$ & MOS$\uparrow$ & SMOS$\uparrow$ \\
        \midrule
        Groundtruth &  & 1.9 & 0.93 & 4.5 & 3.96 \\
        % \midrule
        VALL-E & EnCodec& 7.9 & 0.75 & 3.08 & 3.31 \\
        USLM & \method & \textbf{6.5} & \textbf{0.84} & \textbf{3.63} & \textbf{3.45} \\
        \bottomrule
    \end{tabular}}
    \caption{Results of zero-shot TTS
    }
    \label{tab:zstts}
\end{table*}